\scanpage{119}
% archive sheet number: 118
surjectives, sont-elles en \uline{fait} \uline{bijectives} ?? On \unsure{partira} alors du ss-groupe \inserted{inv} $N$ de $\Gamma$, \unsure{?|\&} et de la carte homogène correspondante $\Gamma/N\simeq\Psi_{\unsure{**|\bullet *}}(P)$

On voudrait que ce soit l'une des \inserted{5} cartes correspondant aux polyèdres \uline{réguliers} connus --- en particulier, on voudrait que ce soit \uline{orientable} ! Et il faudrait, pour une carte combinatoire donnée (d'un de ces cinq types ?) déterminer les polyèdres \uline{réguliers} \uline{épinglés} par ce $\Phi$ --- ce qui est un Pb modulaire \ldots.

\uline{NB} Si la réponse à la question précédente (26) était positive (et sinon, il faudrait renforcer la définition pour qu'elle le devienne) on récupérerait $P$ par la donnée de
\[
 \Gamma\xrightarrow{\varphi}\mathrm{Aut}(X)\tag*{(27)}
\]
et de
\[
 r\equiv(s_0,a_0,f_0)\in\mathrm{Rep}(X)\tag*{(28)}
\]
satisfaisant
\[
 \deleted{\unsure{\sigma_0(a_0)}}\quad
 \left\{\begin{aligned}
 \sigma_0a_0&=a_0,\quad\sigma_0f_0=f_0\\
 \sigma_1s_0&=s_0,\quad\sigma_1f_0=f_0\\
 \sigma_2a_0&=a_0,\quad\sigma_2(\unsure{f_0|s_0})=f_0
 \end{aligned}\right.\tag*{(29)}
\]
en prenant \unsure{pr|pour} $S,A,F$ l'ens des transformés de $s_0,a_0,f_0$ sous $\Gamma$.
